%======================================================================
\section{Revisiones r3 a r6}\label{sec:v8r6}
Estas revisiones salieron de comparar la implementación con Virtual-Memory Powersort
(Moltmann, Nakajima y Wild, ESA 2026) con la misma memoria. Las tres primeras (r3: giros
estimados por muestreo; r4: mezcla por páginas; r5: runs paginados) solo eligen o reorganizan
mezclas estables ya cubiertas; sus argumentos están en \texttt{docs/LEMAS\_v8r.md}. Esta
sección cubre lo que r6 agrega.

\subsection{Tope de particiones malas}
\begin{observacion}\label{obs:tope}
Si el presupuesto inicial de \textsc{QuickEstable} pasa de $\lfloor\lg n\rfloor+1$ a
$\min(\lfloor\lg n\rfloor+1,4)$, la corrección y la terminación no cambian.
\end{observacion}
\begin{proof}
El presupuesto solo decide cuándo un rango pasa al respaldo por mezclas, que ordena de forma
estable cualquier entrada. Sigue siendo un natural que decrece en cada paso desparejo
(incluida la separación de iguales), así que la cantidad de pasos desparejos por camino es
finita, y el resto del argumento de profundidad no depende de su valor inicial.
\end{proof}

\subsection{Una pasada de distribución por páginas}
Sea \seg{A}{0}{m} el rango, $k\ge2$ grupos, páginas de $P\ge1$ elementos y un buffer con $k$
páginas de salida $O_0,\dots,O_{k-1}$, inicialmente vacías. El grupo de $x$ es
$g(x)=|\{j: s_j\le x\}|$ para separadores $s_1\le\dots\le s_{k-1}$, de modo que $x\le y$
implica $g(x)\le g(y)$ y los iguales tienen el mismo grupo. La pasada lee $A[r]$ para
$r=0,\dots,m-1$, lo agrega al final de $O_{g(A[r])}$, y si esa página llega a $P$ elementos
la copia a la ranura $w$, es decir a \seg{A}{wP}{(w+1)P}, anota $G[w]=g$, vacía la página y
hace $w\gets w+1$.

\begin{lema}[ranura libre]\label{lem:ranura}
Cada copia a la ranura $w$ escribe solo posiciones ya leídas y que no contienen elementos
copiados antes.
\end{lema}
\begin{proof}
Antes de leer $A[r]$ vale el invariante $r=wP+\sum_g|O_g|$, porque cada elemento leído está en
exactamente una página de salida o en una ranura copiada. Si al agregar $A[r]$ una página llega
a $P$, entonces $r+1\ge wP+P$, así que \seg{A}{wP}{(w+1)P} está contenido en las posiciones ya
leídas \seg{A}{0}{r+1}. Las ranuras $0,\dots,w-1$ son las únicas escritas antes, y son
disjuntas de la ranura $w$.
\end{proof}

\begin{lema}[orden]\label{lem:orden}
Para cada grupo $g$, la sucesión de sus elementos en orden de entrada es la concatenación de
las ranuras $i$ con $G[i]=g$ en orden creciente de $i$, seguida de $O_g$ al terminar.
\end{lema}
\begin{proof}
$O_g$ recibe los elementos del grupo $g$ en orden de lectura y se vacía entera en cada copia;
las copias usan ranuras con índice estrictamente creciente.
\end{proof}

\begin{lema}[espacio final]\label{lem:espacio}
Al terminar, \seg{A}{wP}{m} mide $\sum_g|O_g|$ y no contiene elementos pendientes.
\end{lema}
\begin{proof}
Es el invariante del Lema~\ref{lem:ranura} con $r=m$.
\end{proof}

\begin{teorema}\label{thm:pasada}
Después de (1) copiar $O_0,\dots,O_{k-1}$ en ese orden a \seg{A}{wP}{m}, (2) permutar las
ranuras con destino $F_{G[i]}+|\{i'<i: G[i']=G[i]\}|$, donde $F_g$ es la cantidad de ranuras
de grupos menores que $g$, y (3) para $g=k-1,\dots,0$, poner $O_g$ en
$[F_{g+1}P+D_g,\ F_{g+1}P+D_{g+1})$ y correr las ranuras del grupo $g$ $D_g$ lugares a la
derecha, con $D_g=\sum_{h<g}|O_h|$, el rango contiene los elementos ordenados de forma estable
por $g$: primero todos los del grupo 0 en orden de entrada, después los del 1, etcétera.
\end{teorema}
\begin{proof}
(1) es posible por el Lema~\ref{lem:espacio}. El destino de (2) es una biyección de
$\{0,\dots,w-1\}$ en sí mismo que agrupa por $G$ y conserva el orden relativo dentro de cada
grupo; seguir sus ciclos con una página de buffer la realiza moviendo cada página una vez. En
(3), el grupo $g$ termina en $[F_gP+D_g,\ F_{g+1}P+D_{g+1})$; como $D_g\le D_{g+1}$ y se procede
de derecha a izquierda, ninguna escritura pisa una página o parcial que todavía no se movió.
Por el Lema~\ref{lem:orden}, cada grupo queda en orden de entrada; los grupos quedan en orden
creciente de $g$, y $g$ es monótono en la clave.
\end{proof}

\begin{observacion}
Comparaciones: $m\lceil\lg k\rceil$ (el grupo se obtiene bajando un árbol implícito de los
separadores). Movimientos: $2m+(M+C)P+bP+2\sum_g|O_g|\le 4{,}5m+O(kP)$, con $M$ páginas
fuera de lugar, $C\le M/2$ ciclos de largo al menos 2 y $b$ páginas corridas. Verificado en
forma exhaustiva para claves en $\{0,\dots,k-1\}$, $k=2..4$, $P=1..3$ y $m\le9$
(1\,140\,207 casos).
\end{observacion}

\subsection{El samplesort}
\textsc{SampleSortPáginas} aplica la pasada con separadores tomados de una muestra y recurre en
cada grupo; un rango que entra en el buffer, un grupo que se queda con todo el rango, un camino
con más de $2\log_k n+2$ pasadas o con dos pasadas seguidas cuyo grupo mayor supera la mitad
van a \textsc{QuickEstable}.

\begin{corolario}
\textsc{SampleSortPáginas} ordena de forma estable y termina.
\end{corolario}
\begin{proof}
Por inducción en la profundidad: cada pasada deja los grupos contiguos, ordenados entre sí y
estables dentro (Teorema~\ref{thm:pasada}); ordenar cada grupo de forma estable da el orden
estable del rango. Cada llamada recursiva es sobre un rango estrictamente menor (si no, va al
caso base) y la profundidad está acotada; los casos base son \textsc{QuickEstable}, ya
demostrado.
\end{proof}

El mecanismo de reusar bloques consumidos y anotar su grupo es de Radsort (Clausecker y
Schintke, 2026), para radix LSD; agrupar en bloques de un tipo y restituir su orden por ciclos
está en Logsort (aphitorite); el árbol de separadores sin saltos, en IPS$^4$o.
